\begin{abstract}
Managing visitor flows is a key challenge for smart cities and smart tourism destinations, particularly in small historic towns where tourism pressure can affect urban sustainability and public-space management. While Artificial Intelligence (AI) techniques have been widely applied to mobility forecasting in large metropolitan areas, limited evidence is available for small destinations characterized by sparse sensing infrastructures and heterogeneous data ecosystems.
This paper presents an AI-driven multi-source forecasting framework for visitor-flow analytics in the historic borough of Dozza, Italy. The proposed framework integrates pedestrian camera counts,  mobile-network indicators, weather observations and event-related information into a unified  data platform supporting multi-horizon forecasting. The study evaluates multiple machine learning models for people forecasting in Dozza.

Results show that  tree-based ensemble models achieve substantial reductions in forecasting error compared with the  baseline. Conversely, demographic indicators derived from mobile-network data remain strongly driven by temporal inertia. The findings demonstrate the value of AI-enabled urban analytics for supporting data-driven decision making in small smart tourism destinations and provide a replicable framework for sustainable visitor management in historic urban environments.
%We evaluate three target families: pedestrian entrances and exits, mobile-network nationality indicators, and age-band presence indicators. 

%The experimental design uses chronological train/test splits, rolling top-k feature validation, leakage-aware feature selection, strong persistence baselines, generalized linear models, tree ensembles, gradient boosting models, a two-stage ridge model, ablation, permutation importance, and SHAP summaries. 



\keywords{Smart Cities \and Artificial Intelligence \and Smart Tourism \and Urban Analytics \and  Flow Forecasting \and Mobile Network Data}
\end{abstract}
